\documentclass[12pt]{article}
\usepackage[utf8]{inputenc}
\usepackage[T1]{fontenc}
\usepackage{hyperref}
\title{Sauraseni Prakrit Language Resources for K-12 Curriculum Development}
\author{Compiled by Claude Sonnet 5.5 (claude-sonnet-5-5) for Andr\'as Kornai}
\date{}
\begin{document}
\maketitle

\section{Introduction}
Sauraseni (\'Saurasen\={\i}) Prakrit (ISO 639-3 \texttt{psu}, Glottocode saur1252) is a historical Middle Indo-Aryan literary language with no native speakers. It appears in Sanskrit dramas (speech of the Vidu\d{s}aka and women characters) and in Digambara Jain canonical texts such as the \d{S}a\d{t}kha\d{n}\d{d}\=agama. Search found only scholarly articles; no school materials exist in the results.

\section{Primary Educational Resources}
No primers, textbooks or school programmes were found.

\section{Linguistic and Reference Materials}
\begin{itemize}
\item Jain Journal, April 1996 (Banerjee), testing Digambara Jain texts against the rules of the Prakrit grammarians (Vararuci, Hemacandra and others). \url{https://jainqq.org/pagetext/Jain_Journal_1996_04/520122/10}
\item Further pages of the same issue discussing Sauraseni features. \url{https://jainqq.org/explore/520122/5}
\item Another scanned Jain text page returned for Sauraseni. \url{https://www.jainqq.org/explore/006968/168}
\end{itemize}

\section{Religious and Community Resources}
The Digambara Jain canon (e.g.\ \d{S}a\d{t}kha\d{n}\d{d}\=agama) is the main textual corpus, as described in the article above; no edition links were verified in this search.

\section{Digital and Audio Resources}
None found. (An LLM-generated encyclopedia entry appeared in results and is deliberately not cited as unreliable.)

\section{Curriculum Development Context}
As a classical literary language, any curriculum would be text-based (reading selected dramatic and Jain passages with Sanskrit/English glosses). Grammarian sources such as Hemacandra's rules (e.g.\ intervocalic \emph{t} $\to$ \emph{d}, as in \emph{tata\d{h}} $\to$ \emph{tado}) give a starting point; teaching editions would need to be located separately.

\end{document}